\setcounter{figure}{0}
\setcounter{table}{0}

\paragraph{Baseline pricing and monitoring's cream skimming effect}

The monitoring firm may wish to increase prices for unmonitored drivers both to encourage monitoring adoption and to reflect a cream skimming effect---advantageous selection into monitoring may draw safer drivers out of the unmonitored pool, raising the latter's average risk. To test this, we exploit the introduction of monitoring in the focal state using a regression discontinuity design. We compare prices and average costs in the unmonitored pool a year before and after monitoring began, controlling for seasonality and observable driver and coverage characteristics. To avoid contamination from attrition, we restrict attention to the first period ($t = 0$). Formally, we estimate:
\begin{align}
dep.\text{ }var._{i} & =\alpha+\gamma Qtr_{i}+\kappa\mathbf{1}_{post,i}+\theta\cdot Qtr_{i}\times\mathbf{1}_{post,i}+\mathbf{x}_{i}^{\prime}\beta+\xi_{y,i}+\epsilon_{i}\label{eq:cs_claim}
\end{align}
where $Qtr_i$ denotes the driver's arrival quarter and $post_i$ indicates whether entry occurred after monitoring was introduced. The monitoring opt-in rate is plotted in \Cref{fig:app_event_study}(a). We examine both price ($p_i$) and claim count ($C_i$) as dependent variables. The coefficient $\theta$ captures the effect of the introduction of the monitoring program on the unmonitored pool.

Estimates for $\hat{\theta}$, reported in \Cref{fig:app_event_study}(b), show no statistically significant price increases and or average cost inflation in the unmonitored pool. Since monitoring represents only a small share of the market, even strong selection into monitoring has limited effect on unmonitored drivers. Moreover, the firm does not pursue customers who initially opt out, making full unraveling unlikely. Finally, because monitoring programs require approval from state commissioners, any baseline price changes are subject to regulatory scrutiny. Taken together, these results suggest that the current regime is largely welfare-neutral for unmonitored drivers.

\begin{figure}[htbp!]
\begin{centering}
\begin{tabular}{cc}
\includeifexists[scale=0.55]{\exhibitpath/appendix/fig_b1a.png} &\hspace{0.05cm} \includeifexists[scale=0.55]{\exhibitpath/appendix/fig_b1b.png}\tabularnewline
(a) &\hspace{0.05cm} (b)
\end{tabular}
\caption{Monitoring Opt-In Rate and Price/Claim Effect Around Introduction \label{fig:app_event_study}}
\par\end{centering}
\vspace{0.1cm}
\begin{minipage}{\textwidth} % choose width suitably\vspace{0.25cm}
    {\footnotesize \emph{Notes: }(a) plots the progression of monthly monitoring finish rate around the introduction of monitoring in the focal state. The monthly finish rate was below $0.1\%$ in all months before monitoring introduction. The reason why it is not exactly zero is due to small-scale trials.
    (b)  reports regression-discontinuity estimate $\theta$ of \eqref{cs_claim}, where the horizontal axis distinguishes dependent variable used; while different colors and positions represent different specifications of control variables ($x_{it}$). The grey (left-most) series represents estimates from regressions with the full set of $x_{it}$; the orange (middle) one includes a full set of observables, including flexible controls for trend and seasonality.  These effects are translated in percentage terms by dividing the average of the dependent variable in the period immediately before monitoring introduction. We look at only first period outcomes, and include all \emph{unmonitored} drivers arriving at the monitoring firm in the focal state a year before or after monitoring introduction. The running variable is quarter since monitoring introduction.
    \par}
\end{minipage}
\end{figure}

\paragraph{Ex-post monitoring pricing and rate revisions}

For each monitored driver, the firm processes the driving data and derive a one-dimensional monitoring score.
Ex-post monitoring pricing maps the driver's monitoring score to a persistent discount or surcharge.

Our focal-state panel covers three baseline pricing regimes and three monitoring pricing regimes. As mentioned above, changes in monitoring pricing across monitoring regimes resulted from both the generation of monitoring scores (\Cref{fig:score-discount-by-regime}a) and how the firm set monitoring discounts corresponding to monitoring tiers (\Cref{fig:score-discount-by-regime}b). In terms of the former, we show that monitoring scores increase with our estimated latent accident risk, but much more so for those produced by the second and third monitoring regimes (\Cref{fig:score-vs-risk-by-regime}). Such increased precision across monitoring regimes was at least partially priced in: across monitoring regimes, monitoring has always increased the already positive relationship between renewal prices and latent accident risk that previously relied only on characteristics-based risk rating, and such increase is further strengthened in latter monitoring regimes (\Cref{fig:R-vs-risk-by-regime}).

\begin{figure}[htbp!]
\begin{centering}
\begin{tabular}{cc}
\includeifexists[scale=0.6]{\exhibitpath/appendix/fig_b2a.png} &\hspace{0.05cm} \includeifexists[scale=0.6]{\exhibitpath/appendix/fig_b2b.png}\tabularnewline
(a) &\hspace{0.05cm} (b)
\end{tabular}
\caption{Monitoring Score: Distribution and Discount Mapping by Monitoring Regimes \label{fig:score-discount-by-regime}}
\par\end{centering}
\vspace{0.1cm}
\begin{minipage}{\textwidth} % choose width suitably\vspace{0.25cm}
    {\footnotesize \emph{Notes: }(a) This graph plots the density of monitoring scores across the three monitoring regimes. Our counterfactual simulations are based on the third monitoring regime. (b) This graph is a binned scatter plot of the monitoring discounts given to monitored drivers against their monitoring score. The Y axis is monitoring renewal factor, where 0.8 implies a 20\% monitoring discount and 1.1 means a 10\% surcharge. Our counterfactual simulations are based on the third monitoring regime.
    \par}
\end{minipage}
\end{figure}

\begin{figure}[H]
    \centering
    \includeifexists[width=0.8\textwidth]{\exhibitpath/appendix/fig_b3.png}
    \caption{Monitoring Score vs. Accident Risk by Monitoring Regimes}
    \label{fig:score-vs-risk-by-regime}
\vspace{0.1cm}
\begin{minipage}{0.95\textwidth} % choose width suitably
    {\footnotesize \emph{Notes: }The X-axis represents the estimated Poisson risk arrival rate while the Y-axis plots monitoring score, data or predicted by our model, by monitoring regimes. Figure \ref{fig:risk_rating_selection}(b) corresponds to monitoring regime 3 in the right panel.\par}
\end{minipage}
\end{figure}

\begin{figure}[H]
    \centering
    \includeifexists[width=0.9\textwidth]{\exhibitpath/appendix/fig_b4.png}
    \caption{Renewal Price vs. Accident Risk by Pricing and Monitoring Regimes}
    \label{fig:R-vs-risk-by-regime}
\vspace{0.1cm}
\begin{minipage}{0.95\textwidth} % choose width suitably
    {\footnotesize \emph{Notes: }The X-axis represents the estimated Poisson risk arrival rate while the Y-axis plots the average first-renewal price for consumers in various risk-bins. We focus on the standard \$50,000 limit plan.  Each panel corresponds, chronologically, to a pricing regime for the baseline (unmonitored) consumer pool in grey round dots. The colored triangle dots plot the monitoring pricing regimes for consumers that opt into monitoring. Figure \ref{fig:risk_rating_selection}(a) corresponds to right-most panel. \par}
\end{minipage}
\end{figure}
